\begin{figure}[htbp]
\centering
\begin{tikzpicture}[x=1cm,y=1cm]
\begin{scope}[shift={(-2.3,0)}]
\draw[AnalysisBlue,very thick] (0,0) circle(1.25);
\fill (0,0) circle(1.2pt);
\draw[->] (0,0)--(1.2,0) node[midway,above]{\(ct\)};
\node[align=center] at (0,-1.9){三维：球面支集\\[2pt]图示为平面截面};
\end{scope}
\begin{scope}[shift={(2.3,0)}]
\fill[AnalysisBlue!16] (0,0) circle(1.25);
\draw[AnalysisBlue] (0,0) circle(1.25);
\fill (0,0) circle(1.2pt);
\draw[->] (0,0)--(1.2,0) node[midway,above]{\(ct\)};
\node[align=center] at (0,-1.9){二维：闭圆盘支集\\[2pt]内部具有非零核};
\end{scope}
\end{tikzpicture}

\smallskip
\begin{tikzpicture}[x=1cm,y=1cm]
\begin{scope}[shift={(-2.3,0)}]
\draw[AnalysisBlue,thick] (0,0) circle(1);
\draw[->] (-1.3,0)--(1.5,0) node[right]{\(z\)};
\draw[->] (0,-1.2)--(0,1.3) node[above]{\(\theta_3\)};
\draw[->,dashed] (.55,{sqrt(1-.55*.55)})--(.55,.04);
\draw[->,dashed] (.55,{-sqrt(1-.55*.55)})--(.55,-.04);
\fill (.55,{sqrt(1-.55*.55)}) circle(1.2pt);
\fill (.55,{-sqrt(1-.55*.55)}) circle(1.2pt);
\node at (0,-1.4){同一圆盘位置};
\end{scope}
\node[align=left] at (2.5,.25){上下半球各贡献一次\\
\(\mathrm d\sigma=(1-|z|^2)^{-1/2}\mathrm dz\)\\
靠近边缘，面积权重增大};
\end{tikzpicture}
\caption[二维与三维初速度核的支集]{二维与三维初速度核的支集示意，不按比例。上排左图只画球面的平面截线；右图为真实二维空间。下排为单位球面的过心截面，\((z,\theta_3)\) 表示球面方向坐标，\(z\) 为二维投影坐标，图中只画其中一个方向。垂直投影把上下两点送到同一圆盘位置；投影圆盘仅显示为一条直径。球面靠赤道趋于竖直，投影面积变小，故球面面积相对于投影面积的权重增大。比较的是固定正时间的核，不是任意初值的实际非零集合。}
\label{b2:fig:11-kernel-support}
\end{figure}
